\clearpage
\onecolumn
\section{VLM Prompt and Coordinate Protocol}
\label{sec:supp-vlm-protocol}

The VLM protocol separates answer preservation from point grounding. Each
render pair is evaluated with immutable source scenes, scratch noMDL
conversion, matched cameras, target projection, blind visual QA, and a fixed
prompt format. The prompt asks for a target-category answer and a
normalized-1000 point. The parser then evaluates whether the answer is scorable
and whether the point lies inside the projected target box.

\begin{center}
\includegraphics[width=\textwidth,height=0.76\textheight,keepaspectratio]{figures/fig_supplement_vlm_coordinate_protocol_atlas.png}
{\captionsetup{hypcap=false}
\captionof{figure}{VLM coordinate protocol. The panels use real
GRScenes target-view renders to separate three visual quantities that can be
confused in table-only form: the target box, the raw image center,
and the normalized-1000 target-box center. The route diagram summarizes the path
from render evidence to parser and scoring steps.}}
\label{fig:supp-vlm-coordinate-protocol-atlas}
\end{center}

\clearpage

\begin{table*}[t]
\centering
\caption{Coordinate-frame ablation for GRScenes VLM pilot grounding. All rows were prompted for normalized-1000 coordinates; raw point columns are diagnostic rescoring in image-pixel space and are not the final benchmark metric. Values are original/converted hit counts unless marked as pair agreement.}
\label{tab:grscenes_vlm_coordinate_ablation}
\scriptsize
\resizebox{\textwidth}{!}{%
\begin{tabular}{lllccccp{0.24\textwidth}}
\toprule
\textbf{Split} & \textbf{Model} & \textbf{Answer O/C} & \textbf{Raw point O/C} & \textbf{Norm-1000 point O/C} & \textbf{Raw pair} & \textbf{Norm pair} & \textbf{Interpretation} \\
\midrule
clean & Gemma4 & 15/15 / 15/15 & 0/15 / 0/15 & 8/15 / 6/15 & 15/15 same; 0/15 both-hit & 11/15 same; 5/15 both-hit & normalized-1000 scoring is the intended protocol \\
clean & Qwen2.5-VL & 8/11 / 9/12 & 5/14 / 5/15 & 0/14 / 0/15 & 9/14 same; 2/14 both-hit & 14/14 same; 0/14 both-hit & raw diagnostic scores higher; coordinate semantics unresolved \\
zoom-stress & Gemma4 & 14/14 / 14/14 & 1/14 / 1/14 & 11/14 / 13/14 & 14/14 same; 1/14 both-hit & 12/14 same; 11/14 both-hit & normalized-1000 scoring is the intended protocol \\
zoom-stress & Qwen2.5-VL & 12/14 / 10/12 & 9/14 / 6/13 & 3/14 / 3/13 & 8/13 same; 5/13 both-hit & 9/13 same; 1/13 both-hit & raw and normalized diagnostics disagree; coordinate semantics unresolved \\
\bottomrule
\end{tabular}%
}
\end{table*}


\begin{center}
\includegraphics[width=\textwidth,height=0.50\textheight,keepaspectratio]{figures/fig_supplement_vlm_coordinate_table_companion.png}
{\captionsetup{hypcap=false}
\captionof{figure}{VLM coordinate table companion. Original/noMDL
render pairs show why raw-image center, target-box center, and normalized-1000
coordinates must be separated in the coordinate ablation table. The diagram
explains the frame transformation.}}
\label{fig:supp-vlm-coordinate-table-companion}
\end{center}

The coordinate ablation explains the main paper's choice of normalized-1000
point hits as the primary metric. A target-centered camera can make the raw
image center a strong baseline even when a model ignores the requested
normalized coordinate convention. The normalized-1000 protocol is stricter
because it tests whether the model follows the requested coordinate format.

\clearpage

\begin{table}[t]
\centering
\caption{Coordinate-only baselines for the expanded30 VLM grounding stress set. These rows are sanity baselines for point scoring, not VLM evidence.}
\label{tab:vlm_coordinate_baselines}
\scriptsize
\resizebox{\linewidth}{!}{%
\begin{tabular}{llllll}
\toprule
\textbf{Baseline} & \textbf{Space} & \textbf{Raw O/C} & \textbf{Norm O/C} & \textbf{Raw agree} & \textbf{Norm agree} \\
\midrule
random\_seeded\_pixel & pixel & 5/30 / 2/30 & 8/30 / 3/30 & 23/30 & 21/30 \\
image\_center\_pixel & pixel & 30/30 / 30/30 & 0/30 / 0/30 & 30/30 & 30/30 \\
bbox\_center\_pixel & pixel & 30/30 / 30/30 & 0/30 / 0/30 & 30/30 & 30/30 \\
bbox\_center\_normalized\_1000 & normalized\_1000 & 0/30 / 0/30 & 30/30 / 30/30 & 30/30 & 30/30 \\
\bottomrule
\end{tabular}%
}
\end{table}


\begin{center}
\includegraphics[width=\textwidth,height=0.50\textheight,keepaspectratio]{figures/fig_supplement_vlm_coordinate_baseline_sanity_companion.png}
{\captionsetup{hypcap=false}
\captionof{figure}{Coordinate-baseline sanity companion. Real target-view
renders show why the deterministic rows in Table S3 are scorer sanity checks:
image-center, bbox-center, and normalized-1000 bbox-center baselines test
different coordinate assumptions. The baseline audit strip adds
original/noMDL target-view pairs for the same scorer checks. The companion
summarizes the reading rule.}}
\label{fig:supp-vlm-coordinate-baseline-sanity-companion}
\end{center}

The deterministic baselines are scorer sanity checks rather than model
predictions. If an oracle box center scores 30/30 under the
normalized format, then VLM failures point to model output behavior rather than
an impossible box or an unavailable scoring frame. If the raw image center also scores 30/30,
then raw-image point hits must be interpreted carefully because the camera
construction itself can make raw coordinates easy.

\subsection{Parser and Scorable Rows}

The parser is conservative by construction. A row is scorable only when the model
produces an answer and a parseable point in the expected form. Scorable
denominators are reported explicitly because a model can fail by returning a
plausible answer with no usable point, by returning a raw pixel coordinate under
the normalized prompt, or by mixing coordinate semantics. These are different
failure modes and are reported separately.

\subsection{Gemma4 and Qwen2.5-VL Interpretation}

Gemma4 is more consistent with normalized-1000 in the frozen stress run, while
Qwen2.5-VL shows a stronger raw-coordinate preference under the same prompt
family. The comparison should be read as backend-specific protocol behavior:
conversion studies must report prompt format, parser coverage, and coordinate
frame together so material deltas are not mistaken for coordinate-convention
effects.
